\section{Conclusion}
\label{sec:conclusion}

Sparse zebrafish fluorescence data require a distinction between event detection, spatial localization, source identification, and biological interpretation. The NeuRev evidence is most coherent when framed as a measurement and identifiability study: temporal differentiation detects burst-associated change; Raw traces contain a strong shared waveform; local subtraction retains event-associated amplitude; and spatial ICA provides a localized, highly stable reconstructed representation despite unstable individual filters. At confirmed centers, amplitude already localizes the bursts nearly perfectly over the complete trace; local coherence adds its clearest value earlier, during full-field proposal and scarce-budget ranking. Consensus and multiscale persistence are promising only as exploratory cross-burst observability features. Automated displacement, perturbation, circular-shift, and injection tests show that the compact features contain localized and event-aligned measurement structure, while grouped recovery modeling stops short of confirming incremental utility. Frozen measurement-event classes recur across bursts, but current data do not establish their alignment with spatial-ICA morphology or reviewer certainty. By preserving spatial-site identity, respecting positive--unlabeled labels, assigning features to task-specific roles, and reserving precision for exhaustively reviewed data, the framework converts a data-limited detection problem into a sequence of falsifiable scientific questions. Realistic synthetic movies, bounded-field truth, and independent-recording transfer are now higher priorities than further same-recording feature expansion.
